\section*{\S\ 5. Die Zerlegungsidentit\"aten in expliziter Form.}
Spezialisieren wir in (25.) im Falle 4) $\tau$ zu $\sigma+\rho$, so kommt:
\begin{equation}
 \pair{q_\rho A^\sigma}{y^{(1)}\cdots y^{(\rho+\sigma)}}
 =\sum_{i_\rho=i_{\rho-1}+1}^{\tau}
 \eps_{i_\rho}\pair{q_\rho}{y^{(i_1)}\cdots y^{(i_\rho)}}
 \pair{A^\sigma}{y^{(i_{\rho+1})}\cdots y^{(i_{\rho+\sigma})}},
\tag{32}
\end{equation}
wo $\eps_{i_\rho}=(-1)^{\delta_{i_\rho}}$, $\delta_{i_\rho}=\rho(\rho+1)/2+i_1+\cdots+i_\rho$ wird. Dies ist eine allgemeinere Form des Produktsatzes f\"ur Matrizen als (16.) und (17.) und ergibt sich durch Laplacesche Entwicklung der Determinante der Produktenmatrix
\[
 \sum\pm(v^{(1)}y^{(1)})\cdots(v^{(\rho)}y^{(\rho)})(a^{(1)}y^{(\rho+1)})\cdots(a^{(\sigma)}y^{(\rho+\sigma)}).
\]
Der allgemeinere Produktsatz ist also aus den Identit\"aten (20.), bzw. (21.) zusammengesetzt, und das Herausgreifen der speziellen Form (16.), (17.) ist dadurch erkl\"art.
